% Generated draft: requires mathematical and source review.
\documentclass[11pt]{article}
\usepackage[margin=25mm]{geometry}
\usepackage{fontspec}
\setmainfont{Latin Modern Roman}
\usepackage{amsmath,amssymb,mathtools,unicode-math}
\setmathfont{Latin Modern Math}
\usepackage{enumitem,needspace,xurl,hyperref}
\setlength{\parindent}{0pt}
\setlength{\parskip}{5pt}
\setlength{\emergencystretch}{3em}
\newcommand{\N}{\mathbb{N}}
\newcommand{\Z}{\mathbb{Z}}
\newcommand{\Q}{\mathbb{Q}}
\newcommand{\R}{\mathbb{R}}
\newcommand{\C}{\mathbb{C}}
\newcommand{\F}{\mathbb{F}}
\newcommand{\A}{\mathbb{A}}
\newcommand{\PP}{\mathbb{P}}
\newcommand{\E}{\mathbb{E}}
\newcommand{\eps}{\varepsilon}
\newcommand{\dd}{\,\mathrm{d}}
\begin{document}
\section*{UnsolvedMath ID: 11000246}
\textbf{Source problem code:} catalogue entry
\subsection{Definitions and mathematical statement}
Let \(\rho:G\hookrightarrow\Gamma_g\), with \(G\cong\pi_1(S_{2g})\), be the specific injection of the source's Theorem 3.1. The theorem, allowed \(g\), surface notation and ambient group definition are absent. The source identifies \(\widetilde S_{2g}\cong\mathbb H^2\), so its surface-group boundary \(\partial_\infty G\) is the circle at infinity.

Question 3.3 asks whether there exists a continuous map
\[
 F:\partial_\infty G\longrightarrow\mathrm{PML}_0(\Sigma)
\]
which is \(G\)-equivariant:
\[
 F(\gamma\xi)=\rho(\gamma)F(\xi)
       \quad(\gamma\in G,\ \xi\in\partial_\infty G).
\]
The subspace \(\mathrm{PML}_0(\Sigma)\), its subscript, topology and action are not defined. It must not be replaced by all projective measured laminations without source evidence. Only continuity and equivariance are requested; no injectivity, surjectivity or extension of a specified orbit map is added.

The following Section 3.5 gives separability context. In its standard ambient-group formulation, \(\Gamma\) is \(H\)-separable when
\[
 \forall x\in\Gamma\setminus H\ \exists K\le\Gamma:\quad
 [\Gamma:K]<\infty,\ H\le K,\ x\notin K.
\]
The source prints \(g\in G\setminus H\) after naming the ambient group \(\Gamma\); the displayed definition makes that symbol mismatch explicit rather than importing the earlier surface subgroup as the ambient domain. LERF requires separability for every finitely generated \(H\); GFERF requires it for every geometrically finite subgroup in a specified geometric setting.

The source motivates separability by asking whether an essential surface in a finite-volume hyperbolic \(3\)-manifold lifts to an embedding in a finite cover. It says geometrically finite subgroup separability often suffices for applications. It further reports that \(\Gamma_g\) is not LERF for \(g\ge2\), since it contains \(F_2\times F_2\), while \(\Gamma_1\) is LERF because it is virtually free. These are contextual source claims, not the answer to the boundary-map question. The passage ends before the next proposed separability question, whose statement is not reconstructed.
\subsection{Short English statement}
For the particular embedded surface group, ask for a continuous equivariant map from its boundary circle to the source's specified lamination subspace. The embedding theorem and target definition are missing. The subsequent LERF/GFERF discussion is preserved as separate context, including its ambient-symbol issue.
\subsection{Sources}
\begin{itemize}
\item Local UnsolvedMath registry, permanent ID \(11000246\), exact source code \texttt{catalogue entry}, Question 3.3, PDF page 286. No individual proposer or proposed year is supplied. The statement refers to Theorem 3.1 and then enters Section 3.5.
\item Farb, edited volume \emph{Problems on Mapping Class Groups and Related Topics}: \url{https://www.math.uchicago.edu/~farb/papers/mcgbook.pdf}. The volume has 329 explicitly typeset numbered items across chapters, repeated local numbering distinguished by PDF page. The extract includes an A. W. Reid page footer, page 280.
\item Internal references [1] and [36] for low-dimensional subgroup-separability literature; full entries absent. No complete separate citation for Theorem 3.1 is supplied.
\item Extraction pdf-font-structure, accessed 2026-07-29. Source evidence NEEDS\_REVIEW presents an open problem and requires a current primary-literature audit. Registry status open and classification OPEN-TRIAGE; the research summary says only ``Literature status: Truncated.''
\item Dataset attribution: Ulam AI / UnsolvedMath Contributors, CC BY 4.0. Source rights NEEDS\_REVIEW: publicly accessible author/publisher copy, redistribution terms requiring review.
\end{itemize}
\end{document}
